\chapter{Hilbertian spectral theory}
\setcounter{exercisepar}{0}

For every vector space E, we denote by \(1_{E}\) the identity map of E. If E, F and G are vector spaces and if v: \(E \rightarrow F\) and u: \(F \rightarrow G\) are linear maps, the linear map \(u \circ v\) from E to G will sometimes be denoted uv.

Let E be a pre-Hilbert space. We denote by \(\langle x|y\rangle_{E}\) , or simply \(\langle x|y\rangle\) , the inner product on E. The orthogonal of a subset A of E is denoted A°. It is a closed subspace of E; it is reduced to 0 if and only if A is total in E (EVT, V, p. 16, cor. 1). If F ⊂ E is a vector subspace of E, then F°° = \(\overline{F}\) (cf. EVT, V, p. 13). The spectrum of an endomorphism of E is always relative to the Banach algebra \(\mathcal{L}(E)\) .

Let F be a Hilbert space. We denote by \(\mathcal{L}_1(\mathrm{E})\) the space of endomorphisms of finite trace of E (EVT, V, p. 50, def. 8) and \(\mathcal{L}_2(\mathrm{E};\mathrm{F})\) the space of Hilbert-Schmidt maps from E to F (EVT, V, p. 51, def. 9). The map \(u\mapsto \| u\| _2 = (\mathrm{Tr}(u^{*}u))^{1 / 2}\) is a Hilbertian norm on \(\mathcal{L}_2(\mathrm{E};\mathrm{F})\) (EVT, V, p. 52, th. 1, (i)).

The space \(\mathcal{L}^{\mathrm{f}}(\mathrm{E};\mathrm{F})\) of continuous linear maps of finite rank is dense in \(\mathcal{L}^{\mathrm{c}}(\mathrm{E};\mathrm{F})\) (prop. 10 of III, p. 15 and cor. of prop. 11 of III, p. 16). In particular, the space \(\mathcal{L}_2(\mathrm{E};\mathrm{F})\) is dense in \(\mathcal{L}^{\mathrm{c}}(\mathrm{E};\mathrm{F})\) .

The Hilbertian tensor product of E and F (EVT, V, p. 28, def. 1) is denoted E \(\widehat{\otimes}_2\) F.

If X is a locally compact topological space and \(\mu\) a measure on X, we denote by \(\mathcal{L}^{p}(\mathrm{X},\mu)=\mathcal{L}_{\mathbf{C}}^{p}(\mathrm{X},\mu)\) and \(\mathrm{L}^{p}(\mathrm{X},\mu)=\mathrm{L}_{\mathbf{C}}^{p}(\mathrm{X},\mu)\) for \(p\in[1,+\infty]\) .

